\subsection{Characteristic and lattice input}
The index and signature polynomials used below are
\cite[equations (1.2)--(1.4), pp.~3--4]{ffEsoleKang}; the
rank-at-most-four integral classification is
\cite[Theorem 2.2.1, p.~61]{ffPetersSterk}.  In the almost
complex application, the Todd polynomial and the cotangent-twisted
polynomial are the corresponding elliptic-symbol indices, so their
integrality does not require an integrable complex structure.
For $c_1=c_3=0$ the substitution gives
\begin{equation}
\operatorname{Td}=\frac{3c_2^2-c_4}{720},\qquad
\chi^1=\frac{3c_2^2-31c_4}{180},\qquad
\sigma=\frac{3c_2^2+14c_4}{45}.
\label{maps:eq:F.1}
\end{equation}
In particular $\chi^1-4\operatorname{Td}=-c_4/6$, and
$\operatorname{Td}=(3\sigma-c_4)/48$.  For the rank-four middle
group and $c_4=6$ this forces $\sigma=2$, since
$\sigma\equiv2\pmod {16}$ and $|\sigma|\leq4$.
The cited lattice theorem then gives $I_{3,1}$ and the displayed
polynomials give the complete numerical record
\begin{equation}
 H^4(X;\mathbb Z)\cong I_{3,1},\qquad c_4=6,\qquad
 c_2^2=2,\qquad\sigma=2,\qquad\operatorname{Td}=0.
 \label{maps:eq:F.2}
\end{equation}
These substitutions are the form of the known inputs used in
the construction; the longer credited re-derivation remains
[S6\_09, \nolinkurl{maps:src:F.1.2}].

Mandal's connected-sum criterion is used only in its stated
orientation: $\ell\mathbb{HP}^2\#m(-\mathbb{HP}^2)$ has an
orientation-compatible almost complex structure exactly when
$\ell=2m+1$ and $m$ is odd
\cite[Corollary 1.4 and its proof in \S5]{ffMandal}.
Its normalized tangent classes are those of
\cite[Lemma 2.1 and Theorem 2.4, \S2]{ffMandal}.
The credited application retained in the compendium is
[S6\_09, \nolinkurl{maps:src:F.1.3}].
